% AUTO-GENERATED by the update_record tool. DO NOT EDIT.
% verify_paper regenerates this file from record.json and the run outputs
% and fails on any difference, so manual edits always surface.
\noindent\textbf{Hypothesis 1.} RIKYU の推論 API の open-weight モデル 5 種（kimi-k2.6、kimi-k3、deepseek-v4.1-flash、qwen3.8-27b、qwen3.6-35b）は、SciGym 公開コードの Controller で反復上限を論文の 20 から 100 に引き上げて SciGym-small 137 件を解かせたとき、いずれも平均 Simulation Trajectory Error（STE）が先行研究の 20 反復の値以下になる（反復を増やしても悪化しない）。~\cite[s1.p3]{duan-2025-measuring}, \cite[s3.p1, s3.p2, s3.p3, s3.p4, s3.p5, s3.p6]{scigymrikyuopenmodels-2026-scigym}
\begin{enumerate}
\item[\textbf{Claim 1}] kimi-k2.6 で SciGym-small 137 件を最大 100 反復で解かせたときの平均 STE は、先行研究の 20 反復の値 0.2100 以下である。~\cite[s1.p3, s1.p4]{duan-2025-measuring}, \cite[s3.p2, s3.p7]{scigymrikyuopenmodels-2026-scigym}, \cite[s2.p1]{h4duan-2025-scigym}
  \emph{Rationale:} 反復を 5 倍にしても悪化しないことが示せれば、h1 の kimi-k2.6 の部分を担う。
  \emph{Criterion:} \texttt{\detokenize{kimi-k2.6-small}}.\texttt{\detokenize{trajectory_smape}} $-$ 0.21005 $\leq 0$.
  \emph{Prediction:} $[-0.08, 0.02]$ (20 反復では大半の件が上限まで使い切っていたので余地はあるが、文脈が長くなると推論が乱れることもある。改善は小さめと見る).
  \emph{Observed:} pending.
  \emph{Verdict:} pending.
\item[\textbf{Claim 2}] kimi-k3 で SciGym-small 137 件を最大 100 反復で解かせたときの平均 STE は、先行研究の 20 反復の値 0.2251 以下である。~\cite[s1.p3, s1.p4]{duan-2025-measuring}, \cite[s3.p3, s3.p7]{scigymrikyuopenmodels-2026-scigym}, \cite[s2.p1]{h4duan-2025-scigym}
  \emph{Rationale:} c1 と同じ。h1 の kimi-k3 の部分を担う。
  \emph{Criterion:} \texttt{\detokenize{kimi-k3-small}}.\texttt{\detokenize{trajectory_smape}} $-$ 0.225072 $\leq 0$.
  \emph{Prediction:} $[-0.08, 0.02]$ (c1 と同じ).
  \emph{Observed:} pending.
  \emph{Verdict:} pending.
\item[\textbf{Claim 3}] deepseek-v4.1-flash で SciGym-small 137 件を最大 100 反復で解かせたときの平均 STE は、先行研究の 20 反復の値 0.2405 以下である。~\cite[s1.p3, s1.p4]{duan-2025-measuring}, \cite[s3.p4, s3.p7]{scigymrikyuopenmodels-2026-scigym}, \cite[s2.p1]{h4duan-2025-scigym}
  \emph{Rationale:} c1 と同じ。h1 の deepseek-v4.1-flash の部分を担う。
  \emph{Criterion:} \texttt{\detokenize{deepseek-v4.1-flash-small}}.\texttt{\detokenize{trajectory_smape}} $-$ 0.24046 $\leq 0$.
  \emph{Prediction:} $[-0.08, 0.02]$ (c1 と同じ).
  \emph{Observed:} pending.
  \emph{Verdict:} pending.
\item[\textbf{Claim 4}] qwen3.8-27b で SciGym-small 137 件を最大 100 反復で解かせたときの平均 STE は、先行研究の 20 反復の値 0.2479 以下である。~\cite[s1.p3, s1.p4]{duan-2025-measuring}, \cite[s3.p5, s3.p7]{scigymrikyuopenmodels-2026-scigym}, \cite[s2.p1]{h4duan-2025-scigym}
  \emph{Rationale:} c1 と同じ。h1 の qwen3.8-27b の部分を担う。
  \emph{Criterion:} \texttt{\detokenize{qwen3.8-27b-small}}.\texttt{\detokenize{trajectory_smape}} $-$ 0.247934 $\leq 0$.
  \emph{Prediction:} $[-0.08, 0.02]$ (c1 と同じ).
  \emph{Observed:} pending.
  \emph{Verdict:} pending.
\item[\textbf{Claim 5}] qwen3.6-35b で SciGym-small 137 件を最大 100 反復で解かせたときの平均 STE は、先行研究の 20 反復の値 0.4571 以下である。~\cite[s1.p3, s1.p4]{duan-2025-measuring}, \cite[s3.p6, s3.p7]{scigymrikyuopenmodels-2026-scigym}, \cite[s2.p1]{h4duan-2025-scigym}
  \emph{Rationale:} c1 と同じ。h1 の qwen3.6-35b の部分を担う。
  \emph{Criterion:} \texttt{\detokenize{qwen3.6-35b-small}}.\texttt{\detokenize{trajectory_smape}} $-$ 0.457089 $\leq 0$.
  \emph{Prediction:} $[-0.15, 0.02]$ (20 反復での値が他の 4 モデルより高く、反復を増やす余地が大きいと見る).
  \emph{Observed:} pending.
  \emph{Verdict:} pending.
\end{enumerate}
\noindent\emph{Assumed, so that the claims together imply Hypothesis 1:}
\begin{itemize}
\item temperature 1.0 での 1 回実行の平均 STE（137 件平均）の実行間変動が判定に対して十分小さいこと。件ごとの STE の標準偏差はおよそ 0.27、137 件平均の標準誤差はおよそ 0.02 で、差が 0.02 未満のときは結論が揺らぎうる（c1〜c5）
\item 先行研究の 20 反復の値（同じ API、同じ Controller、同じ評価層での 1 回実行）を比較対象とし、本研究では 20 反復を再実行しないこと（c1〜c5）
\item 反復上限 100 はプロンプトでエージェントに伝えられ、エージェントは任意の反復で提出できる。上限を使い切らずに提出した件の割合と、文脈長の上限（5 モデル共通の 262144 トークン）や時間切れで提出に至らなかった件の数は表で報告するが、判定には用いない。提出に至らなかった件は公式規約どおり不完全モデルの採点値で平均に含める（c1〜c5）
\item RIKYU の推論 API が提供する各モデルが、先行研究の時点と同じ重み・同じ serving 設定であること（c1〜c5）
\item aarch64 向けの libroadrunner 2.7.0 と antimony 2.14.0 が、論文環境の libroadrunner 2.8.0 と同じ軌道と評価値を与えること（c1〜c5）
\item 平均 STE が「反復を増やして悪化しない」を代表すること。RMS（modifier あり / なし）と NTS は同じ実行から表として報告するが、判定基準には用いない（c1〜c5）
\end{itemize}
\noindent\textbf{Hypothesis 2.} 同じ 5 モデルは、反復上限 100 でより大きな系からなる SciGym-large 213 件を解かせたときも、いずれも平均 STE が論文 Table 1（SciGym-small）の最良行 Gemini-2.5-Pro の 0.3212 以下になる。~\cite[s1.p2, s1.p5, s1.p6]{duan-2025-measuring}, \cite[s3.p1]{scigymrikyuopenmodels-2026-scigym}
\begin{enumerate}
\item[\textbf{Claim 6}] kimi-k2.6 で SciGym-large 213 件を最大 100 反復で解かせたときの平均 STE は、論文 Table 1 の最良行（Gemini-2.5-Pro）の 0.3212 以下である。~\cite[s1.p2, s1.p3, s1.p4, s1.p5]{duan-2025-measuring}, \cite[s3.p2]{scigymrikyuopenmodels-2026-scigym}, \cite[s2.p1]{h4duan-2025-scigym}
  \emph{Rationale:} small で最良だった kimi-k2.6 が large でも論文の最良行の水準を保てば、h2 の kimi-k2.6 の部分を担う。
  \emph{Criterion:} \texttt{\detokenize{kimi-k2.6-large}}.\texttt{\detokenize{trajectory_smape}} $-$ 0.3212 $\leq 0$.
  \emph{Prediction:} $[-0.1, 0.1]$ (small での 0.2100（s3.p2）に、系が大きくなる分の悪化を見込む。先行値が無く幅を広く取る).
  \emph{Observed:} pending.
  \emph{Verdict:} pending.
\item[\textbf{Claim 7}] kimi-k3 で SciGym-large 213 件を最大 100 反復で解かせたときの平均 STE は、論文 Table 1 の最良行（Gemini-2.5-Pro）の 0.3212 以下である。~\cite[s1.p2, s1.p3, s1.p4, s1.p5]{duan-2025-measuring}, \cite[s3.p3]{scigymrikyuopenmodels-2026-scigym}, \cite[s2.p1]{h4duan-2025-scigym}
  \emph{Rationale:} c6 と同じ。h2 の kimi-k3 の部分を担う。
  \emph{Criterion:} \texttt{\detokenize{kimi-k3-large}}.\texttt{\detokenize{trajectory_smape}} $-$ 0.3212 $\leq 0$.
  \emph{Prediction:} $[-0.1, 0.1]$ (c6 と同じ（small での値は 0.2251、s3.p3）).
  \emph{Observed:} pending.
  \emph{Verdict:} pending.
\item[\textbf{Claim 8}] deepseek-v4.1-flash で SciGym-large 213 件を最大 100 反復で解かせたときの平均 STE は、論文 Table 1 の最良行（Gemini-2.5-Pro）の 0.3212 以下である。~\cite[s1.p2, s1.p3, s1.p4, s1.p5]{duan-2025-measuring}, \cite[s3.p4]{scigymrikyuopenmodels-2026-scigym}, \cite[s2.p1]{h4duan-2025-scigym}
  \emph{Rationale:} c6 と同じ。h2 の deepseek-v4.1-flash の部分を担う。
  \emph{Criterion:} \texttt{\detokenize{deepseek-v4.1-flash-large}}.\texttt{\detokenize{trajectory_smape}} $-$ 0.3212 $\leq 0$.
  \emph{Prediction:} $[-0.1, 0.1]$ (c6 と同じ（small での値は 0.2405、s3.p4）).
  \emph{Observed:} pending.
  \emph{Verdict:} pending.
\item[\textbf{Claim 9}] qwen3.8-27b で SciGym-large 213 件を最大 100 反復で解かせたときの平均 STE は、論文 Table 1 の最良行（Gemini-2.5-Pro）の 0.3212 以下である。~\cite[s1.p2, s1.p3, s1.p4, s1.p5]{duan-2025-measuring}, \cite[s3.p5]{scigymrikyuopenmodels-2026-scigym}, \cite[s2.p1]{h4duan-2025-scigym}
  \emph{Rationale:} c6 と同じ。h2 の qwen3.8-27b の部分を担う。
  \emph{Criterion:} \texttt{\detokenize{qwen3.8-27b-large}}.\texttt{\detokenize{trajectory_smape}} $-$ 0.3212 $\leq 0$.
  \emph{Prediction:} $[-0.1, 0.1]$ (c6 と同じ（small での値は 0.2479、s3.p5）).
  \emph{Observed:} pending.
  \emph{Verdict:} pending.
\item[\textbf{Claim 10}] qwen3.6-35b で SciGym-large 213 件を最大 100 反復で解かせたときの平均 STE は、論文 Table 1 の最良行（Gemini-2.5-Pro）の 0.3212 以下である。~\cite[s1.p2, s1.p3, s1.p4, s1.p5]{duan-2025-measuring}, \cite[s3.p6]{scigymrikyuopenmodels-2026-scigym}, \cite[s2.p1]{h4duan-2025-scigym}
  \emph{Rationale:} c6 と同じ。h2 の qwen3.6-35b の部分を担う。small では 0.4571 と基準を上回っていたので、反復を増やしても届かない可能性が高い。
  \emph{Criterion:} \texttt{\detokenize{qwen3.6-35b-large}}.\texttt{\detokenize{trajectory_smape}} $-$ 0.3212 $\leq 0$.
  \emph{Prediction:} $[0, 0.25]$ (small での 0.4571（s3.p6）が基準を 0.14 上回っていた。反証されると予測する).
  \emph{Observed:} pending.
  \emph{Verdict:} pending.
\end{enumerate}
\noindent\emph{Assumed, so that the claims together imply Hypothesis 2:}
\begin{itemize}
\item large split には論文の公表値が無いので、論文の small での最良値を絶対水準の基準とする。難易度の差は small の表と並べて示すが、判定には用いない（c6〜c10）
\item 213 件平均の実行間変動が判定に対して十分小さいこと。件ごとの標準偏差を 0.27 とすれば平均の標準誤差はおよそ 0.02（c6〜c10）
\item 文脈長（262144）や時間切れで提出に至らなかった件は不完全モデルの採点値で平均に含める。large は SBML が長く文脈超過が増えうるが、その件数を表で報告する（c6〜c10）
\item 評価層 scigym\_large（airas-eval v0.13.0）が公式リリースの large split 213 件を sha256 で固定し、small と同じ公式 Evaluator で採点すること（c6〜c10）
\end{itemize}
